% ==========================================
% BAB VII PENUTUP
% ==========================================
\cleardoublepage
\chapter{PENUTUP}
\label{chap:penutup}

\section{Kesimpulan}

Berdasarkan hasil analisis, perancangan, realisasi sistem, dan evaluasi yang
telah dilakukan, kesimpulan tugas akhir ini adalah sebagai berikut:

\begin{enumerate}
  \item Tugas akhir ini berhasil merumuskan rancangan arsitektur sistem inspeksi
  kargo digital terintegrasi yang memisahkan tanggung jawab sistem ke dalam
  komponen \textit{edge}, layanan API, aplikasi web, dan lapisan penyimpanan.
  Struktur tersebut selaras dengan kebutuhan modularitas, integrasi, dan
  ketertelusuran proses inspeksi di lingkungan pelabuhan.
  \item Rancangan arsitektur yang disusun berhasil menempatkan data inspeksi
  sebagai inti integrasi sistem. Hal ini direalisasikan melalui model data yang
  memusatkan informasi pada entitas kontainer serta mendukung keterkaitan dengan
  data manifes, hasil pemindaian, dan artefak visual pendukung.
  \item Artefak realisasi sistem menunjukkan bahwa rancangan arsitektur yang
  diusulkan dapat diwujudkan. Komponen \textit{edge}, layanan API, aplikasi web,
  dan media penyimpanan dapat diintegrasikan sehingga membentuk alur data yang
  utuh, mulai dari identifikasi kontainer, pengumpulan hasil pemindaian,
  penyimpanan bukti inspeksi, hingga penyajian hasil kepada pengguna.
  \item Hasil evaluasi menunjukkan bahwa sistem yang direalisasikan telah
  mendukung kebutuhan utama tugas akhir, yaitu integrasi antarkomponen,
  konsistensi alur data, pemisahan jalur komunikasi sesuai fungsi, dan
  ketertelusuran proses inspeksi. Dengan demikian, rancangan arsitektur yang
  dihasilkan dapat menjadi landasan pengembangan sistem inspeksi kargo digital
  terintegrasi, dengan catatan bahwa evaluasi pada tugas akhir ini berfokus pada
  validasi arsitektur sistem, bukan pada penilaian menyeluruh terhadap
  \textit{deployment} produksi, dampak bisnis, atau akurasi model deteksi secara
  mendalam.
\end{enumerate}

\section{Saran}

Berdasarkan hasil tugas akhir dan keterbatasan yang telah diidentifikasi,
pengembangan lanjutan dapat diarahkan pada beberapa hal berikut:

\begin{enumerate}
  \item Melakukan evaluasi kuantitatif terhadap performa sistem, seperti latensi
  menyeluruh, kapasitas pemrosesan, dan ketahanan sistem ketika dioperasikan
  pada beban yang lebih tinggi.
  \item Melakukan pengujian lapangan pada lingkungan pelabuhan yang lebih
  representatif agar rancangan arsitektur dapat divalidasi terhadap kondisi
  operasional nyata, termasuk variasi jaringan, kualitas aliran video, dan
  proses kerja harian petugas.
  \item Mengembangkan integrasi yang lebih luas dengan sistem eksternal di
  ekosistem pelabuhan, seperti sistem operasional terminal, sistem komunitas
  pelabuhan, dan sistem kepabeanan, agar hasil inspeksi menjadi bagian dari alur
  informasi logistik yang lebih menyeluruh.
  \item Menyempurnakan mekanisme tata kelola data dan audit, termasuk
  pengelolaan riwayat perubahan, retensi artefak inspeksi, dan penguatan kontrol
  akses agar sistem lebih siap digunakan pada lingkungan yang memiliki tuntutan
  kepatuhan tinggi.
  \item Melakukan evaluasi lebih mendalam terhadap kualitas model deteksi dan
  OCR pada variasi kondisi kontainer, pencahayaan, dan kualitas citra agar
  performa keseluruhan sistem semakin andal pada implementasi skala operasional.
\end{enumerate}
